\scanpage{109}
\mspage{3.2}
% archive sheet number: 108
\uline{Proof}: 1) Assume that $S$ is finitely generated as $R\unsure{'}$-module. Then
$\mathrm{w.}\ \mathrm{hd}_R S<\infty$ implies $\mathrm{hd}_R S<\infty$. Let $0\longrightarrow X_t\longrightarrow\ldots\longrightarrow X_1\longrightarrow X_0\longrightarrow
S\longrightarrow 0$ be an acyclic resolution of $S$ by $R$-free modules of finite rank,
and let $\underline m$ be the maximal ideal of $R$. Then by the additivity of the multiplicity,
we have $\dfrac{e_S(\underline m)}{e_R(\underline m)}=\handwritten{\sum}_i(-1)^i\,e_{X_i}(\underline m)/e_R(\underline m)=\handwritten{\sum}_i(-1)^i[X_i:R]=$
\[
\begin{aligned}
&\handwritten{\sum}_i(-1)^i[X_i\mathbin{\ooalign{\hfil$\mathrm{x}$\hfil\cr\hfil$\handwritten{\bigcirc}$\hfil\cr}}_R\bar R:R]=\handwritten{\sum}_i(-1)^i[\operatorname{Tor}_i^R(S,\bar R):\bar R]=\qquad\unsure{:}\\
&[\bar S:\bar R]\cdot[\handwritten{\sum}_i(-1)^i L_S(\operatorname{Tor}_i^R(S,\bar R))].
\end{aligned}
\]
Therefore $e(S/R)=\dfrac{e_S(\underline m S)}{e_R(\underline m)}=$
\[
\dfrac{1}{[\bar S:\bar R]}\cdot\dfrac{e_S(\underline m)}{e_R(\underline m)}=\handwritten{\sum}_i(-1)^i L_S(\operatorname{Tor}_i^R(S,\bar R)).
\]

2) General case: Since $\bar S$ is finite dimensional over $\bar R$, $\overset{\handwritten{\scriptscriptstyle\wedge}}{S}$ is finitely
generated as $\overset{\handwritten{\scriptscriptstyle\wedge}}{R}$-module. Furthermore, by lemma 3.1, we have $\mathrm{hd}_{\overset{\handwritten{\scriptscriptstyle\wedge}}{R}}\overset{\handwritten{\scriptscriptstyle\wedge}}{S}<\infty$.
Therefore we are brought to the above situation and hence $e(S/R)=e(\overset{\handwritten{\scriptscriptstyle\wedge}}{S}/\overset{\handwritten{\scriptscriptstyle\wedge}}{R})=$
\[
\handwritten{\sum}_i(-1)^i L_{\overset{\handwritten{\scriptscriptstyle\wedge}}{S}}(\operatorname{Tor}_i^{\overset{\handwritten{\scriptscriptstyle\wedge}}{R}}(\overset{\handwritten{\scriptscriptstyle\wedge}}{S},\bar R))=\handwritten{\sum}_i(-1)^i L_S(\operatorname{Tor}_i^R(S,\bar R)).
\]

Corollary 3.3: Let $S$ be a local $R$-algebra over a local ring $R$. If
$\mathrm{w.}\ \mathrm{hd}_S S^f<\infty$, then $r(S/R)$ is an integer, equal to $\handwritten{\sum}_i(-1)^i L_{S^f}(\operatorname{Tor}_i^S(S^f\underline m\,\bar S))=$
\[
\handwritten{\sum}_i(-1)^i L_{S^f}(S^f\mathbin{\ooalign{\hfil$\mathrm{x}$\hfil\cr\hfil$\handwritten{\bigcirc}$\hfil\cr}}_{S^e}\operatorname{Tor}_i^S(S^e,\bar S)).
\]

Lemma 3.4: Let $R\handwritten{\subset} S$ be local domains such that $S$ is finitely
generated as $R$-module. If $R$ is integrally closed and $\bar S$ is a simple
extension of $\bar R$, then $e(S/R)$ is an integer.

\uline{Proof}: Let $\alpha$ be $\underset{\unsure{:}}{\text{an}}$ element in $S$ such that $\bar S=\bar R(\bar\alpha)$ where $\bar\alpha$ is the
image of $\alpha$ under the canonical map $S\longrightarrow\bar S$. The subring $R[\alpha]$ is $R$-free
since $R$ is integrally closed, and is a local $R$-algebra since $S$ is a finite
local $R$-algebra. Therefore $e(R[\alpha]/R)=[R[\alpha]:R]/[\bar R[\bar\alpha]:\bar R]=$
$[R[\alpha]:R]/[\bar S:\bar R]$ is an integer by proposition 3.2. Consequently
$e(S/R)=[S:R]/[\bar S:\bar R]=[S:R[\alpha]]\cdot e(R[\alpha]/R)$ is an integer.
